% BHU PYQ -- Manually transcribed & error-corrected
% Paper: MTB-202 : Statics & Dynamics
% Exam: B.Sc. (Hons.) Semester II Examination, 2023-24

\documentclass[11pt,a4paper]{article}
\usepackage[a4paper,top=2.5cm,bottom=2.5cm,left=2.8cm,right=2.8cm,headheight=14pt]{geometry}
\usepackage{amsmath,amssymb}
\usepackage[T1]{fontenc}\usepackage[utf8]{inputenc}
\usepackage{lmodern,microtype,fancyhdr,enumitem,hyperref,lastpage,parskip}

\hypersetup{colorlinks=true,urlcolor=black,linkcolor=black,
  pdftitle={MTB-202 Statics \& Dynamics -- BHU B.Sc. Sem II 2023-24}}
\pagestyle{fancy}\fancyhf{}
\renewcommand{\headrulewidth}{0.4pt}\renewcommand{\footrulewidth}{0.4pt}
\fancyhead[L]{\small   \textit{Banaras Hindu University}}
\fancyhead[R]{\small   \textit{MTB-202: Statics \& Dynamics}}
\fancyfoot[C]{\small Page  \thepage\ of \pageref{LastPage}}


\newlist{parts}{enumerate}{2}
\setlist[parts,1]{label=(\alph*),leftmargin=*,itemsep=5pt,topsep=3pt}
\setlist[parts,2]{label=(\roman*),leftmargin=*,itemsep=3pt,topsep=2pt}

\newcommand{\pts}[1]{\hfill   \textit{[#1]}}


\begin{document}

\begin{center}
  {\large\bfseries BANARAS HINDU UNIVERSITY}\\[2pt]
  {\normalsize B.Sc.\ (Hons.) --- Semester II Examination, 2023-24}\\[6pt]
  \rule{0.8\linewidth}{0.5pt}\\[6pt]
  {\large\bfseries Paper No.\ MTB-202: Statics \& Dynamics}\\[4pt]
  {\normalsize Time: 3 Hours \hfill Maximum Marks: 70}\\[2pt]
  {\small Ref.\ No.:\ 074904}
\end{center}
\rule{\linewidth}{0.4pt}
\smallskip



\noindent   \textit{Instructions: Answer   \textbf{five} questions including Question~1 which is compulsory. Figures in right-hand margin indicate marks.}
\bigskip




\noindent   \textbf{Question 1.}   \textit{(Compulsory --- 2 marks each.)}

\begin{parts}
  \item Obtain a single force or couple equivalent to four forces with components $(2,3)$, $(3,-1)$, $(-4,-4)$, $(-1,2)$ acting at $(0,1)$, $(1,-1)$, $(2,3)$, $(-1,-2)$ respectively.
  \item Find the work done by the thrust at the ends of a rod when displacement alters its length.
  \item Prove that a system of coplanar forces is in equilibrium if the algebraic sum of moments about each of three non-collinear points is zero.
  \item A particle of mass $m$ moves along a path with linear velocity $v$ at distance $r$ from the origin. Obtain the angular velocity $\omega$.
  \item Which of the following are NOT properties of SHM: (i) acceleration always directed to a fixed line; (ii) constant distance from a fixed point; (iii) constant linear velocity; (iv) motion always in a straight line.
  \item If radial and transverse velocity components are $\lambda r$ and $\mu r^2$ ($\lambda,\mu$ constants), find the polar equation of the path.
  \item State Kepler's third law in mathematical form and prove it.
  \item A particle describes a cardioid $ap^2=r^3$ under a central force. Find the law of force.
  \item Is areal velocity of a particle under central force always proportional to angular momentum per unit mass? Give reasons.
  \item A particle of mass 3 units moves up the inner surface of a smooth vertical circular path of radius 5 with initial velocity 15. Find the reaction at angle $60°$ ($g=9.8$).
  \item A frame has linear motion with acceleration. Is it a Newtonian frame of reference? Give reasons.
\end{parts}
\medskip\rule{\linewidth}{0.2pt}

\medskip


\noindent   \textbf{Question 2.}

\begin{parts}
  \item Show that a system of coplanar forces can be reduced to a single force through a given point and a couple.\pts{6}
  \item A system of forces in the plane of triangle $ABC$ is equivalent to a single force at $A$ along the internal bisector of $\angle BAC$ and a couple $G_1$. If moments about $B$ and $C$ are $G_2$ and $G_3$, show $(b+c)G_1=bG_2+cG_3$, where $b=CA$, $c=AB$.\pts{6}
\end{parts}
\medskip\rule{\linewidth}{0.2pt}

\medskip


\noindent   \textbf{Question 3.}

\begin{parts}
  \item State and prove the principle of virtual work for a system of coplanar forces.\pts{6}
  \item A heavy elastic string of natural length $2\pi a$ is placed round a smooth cone of semi-vertical angle $ \theta$, axis vertical. Weight $W$, modulus $\lambda$. Prove it is in equilibrium when it forms a circle of radius $a\!\left(1+\dfrac{W an \theta}{2\pi\lambda}\right)$.\pts{6}
\end{parts}
\medskip\rule{\linewidth}{0.2pt}

\medskip


\noindent   \textbf{Question 4.}

\begin{parts}
  \item Derive the tangential and normal components of acceleration of a particle moving in a plane curve. Also write the radial and transverse components.\pts{6}
  \item A point moves in a curve so that tangential and normal accelerations are equal and the tangent rotates with constant angular velocity. Show the intrinsic equation is $s=Ae^{ \theta}+B$.\pts{6}
\end{parts}
\medskip\rule{\linewidth}{0.2pt}

\medskip


\noindent   \textbf{Question 5.}

\begin{parts}
  \item Show that in SHM the average speed and average acceleration are obtained by multiplying their maximum values by $ frac{2}{\pi}$.\pts{6}
  \item Discuss the motion of a particle in a frame rotating with constant angular velocity about a fixed frame and obtain the fictitious forces.\pts{6}
\end{parts}
\medskip\rule{\linewidth}{0.2pt}

\medskip


\noindent   \textbf{Question 6.}

\begin{parts}
  \item If the central orbit is $r^3=a^3\cos 3 \theta$ under force toward the pole, find the law of force and velocity at any position.\pts{6}
  \item A particle under central acceleration $\mu/ \text{(distance)}^2$ is projected with velocity $V$ at distance $R$. Show its path is a rectangular hyperbola if angle of projection is $\sin^{-1}\!\left\{\mu/(VR\sqrt{V^2-2\mu/R})\right\}$.\pts{6}
\end{parts}
\medskip\rule{\linewidth}{0.2pt}

\medskip


\noindent   \textbf{Question 7.}

\begin{parts}
  \item A particle moves on a smooth vertical circular wire of radius $a$, projected from the lowest point with velocity $\sqrt{4ga}$. Show the reaction is zero after time $\sqrt{a/g}\,\log(\sqrt{5}+\sqrt{6})$.\pts{6}
  \item Two particles drop from the cusp of a cycloid at interval $t$. Prove they meet after time $2\pi\sqrt{a/g}+t/2$.\pts{6}
\end{parts}

\vfill\rule{\linewidth}{0.4pt}

\begin{center}\small
\href{https://bhu-exam.vercel.app/}{Click here to visit our website}\\[4pt]
\href{https://me-aryan.vercel.app/}{Click here to visit our contributor}\\[4pt]
\href{https://quashan-me.vercel.app/}{Click here to visit our contributor}
\end{center}
\end{document}
